\section{开源之路：天授与 tuixue.online}

\subsection{天授（Tianshou）：两周手搓的 RL 框架}

2020年初，翁家翌做了一个改变他职业轨迹的决定：从零开始写一个强化学习框架。

起因很简单——他写了很多 RL 实验代码，想整合一下让自己跑得更好。他先看了 Ray 下面的 RLlib，花了一个月发现太复杂了——几十万行代码，抽象太多，完全不知道该怎么改。于是他决定\textbf{推倒重来，手搓一个全新的框架}。

\begin{importantbox}{天授的设计哲学：一致性（Consistency）}
翁家翌认为一个好项目最重要的特性是\textbf{一致性}。多人协作的项目容易腐化，因为每个人都不知道对面写了什么，假设无法及时传递，导致代码复制粘贴和膨胀。天授之所以成功，是因为从头到尾由一个人设计和实现，保证了整体的一致性。
\end{importantbox}

天授的第一版只花了\textbf{两周}就完成了。翁家翌解释：如果把抽象搞对，实现一个算法可能不到二十行代码。相比 RLlib 的几十万行，天授的关键在于精简的顶层设计——用户想改什么功能，设计上已经指定好了"只能改这个地方"。

天授做对的核心是\textbf{抓住了用户需求}：当时 RL 研究者需要一个好用、好改、代码短的框架，能直接拿来用，研究一下就知道该改哪里。天授恰好满足了这个需求。

\subsubsection{天授的后续发展与腐化}

天授后来成为一个开源社区维护的项目。翁家翌在入职 OpenAI 后没有时间继续维护，便将维护权转移给了社区。他设定的规则是：只要有一个拍板的人，就可以保持 consistency。

五年后回看，翁家翌承认天授"有一点"腐化了——因为他的 context 和接任者的 context 不完全一致，接任者会重写部分代码，导致整体不那么 consistent 了。但他认为这在长远来看是可以接受的。

这段经历也让他在 2022 年 8 月做出了一个重要决定：逐步停止天授的开发。原因是他进入 OpenAI 后意识到，天授针对的 Atari、MuJoCo 等 toy benchmark 与工业界真正需要的 RL（如 LLM 对齐）有巨大鸿沟。"我应该投入更多的时间到更有意义的事情里面。"

\begin{warningbox}{关于"不发 Paper"的选择}
翁家翌明确表示做天授不是为了发论文："我觉得发 paper 完全没有意义。"他当时已经有了论文、比赛成绩，申请也够用了。他想要的是一个真正的、三位数 star 以上的开源项目——一个按导师评价体系来说"正儿八经"的 GitHub 项目。
\end{warningbox}

\subsection{tuixue.online：疫情中的签证查询系统}

2020年疫情期间，美国领事馆关闭，签证时间高度不确定。翁家翌自己也在等签证去 CMU，于是写了一个签证时间爬虫，开源为免费查询系统 tuixue.online。

这个项目的第一版甚至不需要什么技术——白天手动更新一次，晚上手动更新一次。但即使如此简陋，需求就已经很大了。后来自动化后，累计点击量超过百万次（后来可能达到千万次）。疫情过后领事馆升级了网站，原来的爬虫用不了，翁家翌也没时间重写，项目就完成了它的历史使命。

翁家翌对这个项目的总结是："技术不重要，重要的就是抓住需求。"你甚至不需要高级的技术——哪怕是手动更新，只要满足了用户的痛点需求，就有巨大的价值。

\subsubsection{开源是一种慈善}

翁家翌反复用"慈善"来定义自己的开源项目。主持人注意到，天授和 tuixue.online 其实都不是功利的项目——天授做的时候申请已经结束了，tuixue 也完全免费。翁家翌承认自己有一种"很强烈的内在冲动"——想要创造一些自己觉得有用的东西，然后分享给所有人。

他分享了这种冲动的源头：高三时突然蹦出一个 idea——"如果人生是一场游戏，那么游戏的结算分数就是你死的那个瞬间记得你名字的人的数量。"这个"人生游戏论"一直驱动着他追求 impact，并且一直延续到今天。

主持人追问："你不觉得自己也被这个标准推着走了吗？"翁家翌的回答很有意思："目前还没有。如果发现这种情况可以改——可以改自己的评价标准。"他不是这个标准的"奴隶"——哪怕在 OpenAI 很长一段时间没有新的开源项目，他也不觉得困扰，"OpenAI model 就是最好的（开源项目）"。

\begin{importantbox}{代码即慈善}
翁家翌将天授和 tuixue.online 都定义为\textbf{慈善项目}（non-profit）。他说："做慈善项目让我感觉非常满足。相比钱，impact 更让我满足。"这种"代码即慈善"的理念，源于他高中时突然意识到的一个"人生游戏结算规则"——\textbf{你死的那个瞬间，记得你名字的人越多，分数越高}。
\end{importantbox}

\subsection{本章小结}

天授和 tuixue.online 共享同一个创作模式：\textbf{自己有需求 → 市面上没有好用的工具 → 手搓一个 → 免费开源给所有人}。这两个项目都不功利——天授不是为了发论文，tuixue 不是为了赚钱。但正是这些"非功利"的项目，为翁家翌带来了最大的 impact 和职业机会。技术不重要，\textbf{抓住需求}才重要。


